\documentclass[11pt]{article}
\usepackage[T1]{fontenc}
\usepackage{lmodern,amsmath,amssymb,amsthm}
\usepackage[margin=1in]{geometry}
\usepackage[hidelinks]{hyperref}
\usepackage{microtype}
\DeclareMathSizes{10}{10}{7}{7}
\hypersetup{pdftitle={Free spanning-forest degree can decrease under squaring},
pdfauthor={Alec Kriebel},
pdfsubject={A nonunimodular vertex-transitive graph-square counterexample},
pdfkeywords={free uniform spanning forest, graph square, grandparent graph, 1200024}}
\newtheorem{theorem}{Theorem}
\newtheorem{lemma}[theorem]{Lemma}
\newtheorem{proposition}[theorem]{Proposition}
\theoremstyle{remark}\newtheorem{remark}[theorem]{Remark}
\newcommand{\FUSF}{\operatorname{FUSF}}
\newcommand{\E}{\mathbb E}
\newcommand{\RF}{R^{\mathrm F}}
\newcommand{\energy}{\mathcal E}
\renewcommand{\labelitemi}{\(\bullet\)}
\title{Free spanning-forest degree can decrease under squaring}
\author{Alec Kriebel\\
\small Independent researcher\\
\small ORCID: \href{https://orcid.org/0009-0001-9320-500X}{0009-0001-9320-500X}}
\date{22 October 2026\\\small Preprint, version 1.0.0}
\begin{document}
\maketitle
\begin{abstract}
The expected root degree of the free uniform spanning forest need not be
invariant under squaring an infinite connected locally finite simple
vertex-transitive graph. For the binary grandparent graph \(G\), of degree
\(8\), two explicit free-exhaustion reductions give
\(\E\deg_{\FUSF(G)}(o)=(24-\sqrt2)/7>3.22\).
For its simple square \(G^2\), of degree \(40\), a finite rational-current
certificate proves
\(\E\deg_{\FUSF(G^2)}(o)\le3208526010913748780/10^{18}<3.21\).
All currents, divergences and energies are verified with integer arithmetic;
the passage to the infinite free forest is justified by finite-support
flows and free exhaustion. With the explicitly adopted normalization
\(b_F(X)=\frac12\E\deg_{\FUSF(X)}(o)-1\), this gives
\(b_F(G)-b_F(G^2)>1/200\), a negative answer to the unrestricted
vertex-transitive assertion in Ab\'ert's Question~24.
The graph and electrical/forest principles are established background;
no new general forest mechanism or exhaustive priority claim is made.
\end{abstract}

\section{The question and the precise convention}
For a simple graph \(X\), its simple square \(X^2\) has the same vertices,
with distinct vertices adjacent when their \(X\)-distance is at most two.
Every edge below has unit conductance.
The free uniform spanning forest, abbreviated FUSF, is the weak limit of
uniform spanning trees on finite connected induced subgraphs exhausting
\(X\), with no boundary vertices identified \cite[\S5]{BLPS}.
For vertex-transitive \(X\), define
\begin{equation}\label{bf}
 b_F(X)=\tfrac12\,\E\deg_{\FUSF(X)}(o)-1.
\end{equation}
The choice of \(o\) does not matter.

Ab\'ert's \emph{Some questions}, dated 2 November 2010, printed p.~4,
Question~24 \cite{Abert}, asks square invariance for the first \(L^2\)
Betti number obtained from expected free-forest degree for an
\emph{arbitrary} vertex-transitive graph. The original does not print a
normalization formula. We state \eqref{bf} explicitly: it is the usual
forest-derived normalization matching
\(\E\deg_{\FUSF}=2+2\beta^{(2)}_1\) on Cayley graphs
\cite[p.~1]{PT}. Our directly proved assertion is unequal expected
degrees, and the answer about \(b_F\) uses \eqref{bf}.
We do not identify a different nonunimodular \(L^2\) invariant with it.
No unimodularity hypothesis is present in the inspected question;
its adjacent comment already affirms the Cayley case.

\begin{samepage}
\begin{theorem}\label{main}
Let \(T\) be the \(3\)-regular tree with a distinguished end, let \(p(v)\)
be the neighbor toward that end, and let \(G\) have the unordered edges
\(\{v,p(v)\}\) and \(\{v,p^2(v)\}\).
Then \(G\) is an infinite connected locally finite simple vertex-transitive
graph of degree \(8\), with nonunimodular full automorphism group.
Its square \(H=G^2\) is vertex-transitive of degree \(40\), and
\begin{align}
 \E\deg_{\FUSF(G)}(o)&=\frac{24-\sqrt2}{7}>\frac{161}{50},\label{gvalue}\\
 \E\deg_{\FUSF(H)}(o)&\le
 \frac{3208526010913748780}{10^{18}}<\frac{321}{100}.\label{hbound}
\end{align}
Consequently \(b_F(G)-b_F(H)>1/200\).
\end{theorem}
\end{samepage}

The grandparent family is prior; Tang studies its FUSF connectivity and
branching \cite[\S5.2]{Tang}. The contribution here is the specific
expected-degree comparison, using the exact original value and a finite
certificate for the square. The known Cayley case and unimodular
framework are not refuted. We make no assertion about a new general
electrical method.

\section{The graph and the electrical facts used}
The end stabilizer in \(\operatorname{Aut}(T)\) acts transitively and
preserves \(G\), proving transitivity; \(G\) contains \(T\).
At any vertex there are three incident tree edges and five incident
grandparent edges, with no repetitions.

We also verify nonunimodularity for the \emph{full} group.
A tree edge \(x\,p(x)\) lies in exactly three triangles of \(G\), through
\(p^2(x)\) and the two children of \(x\).
A grandparent edge \(x\,p^2(x)\) lies in exactly one, through \(p(x)\).
Thus \(G\) intrinsically determines the tree edges.
On the three recovered tree-neighbors of \(x\), the induced \(G\)-graph
is \(K_{1,2}\), whose unique center is \(p(x)\).
Every automorphism therefore preserves \(p\) and the distinguished end.
Conversely an end-preserving tree automorphism preserves \(G\).
The stabilizer of \(x\) fixes \(p(x)\), while the stabilizer of \(p(x)\)
has an orbit of size two on its children, including \(x\).
The paired suborbit sizes \(1,2\) violate the usual unimodularity
criterion \cite[\S2.3]{Tang}. A nonunimodular transitive subgroup alone
would not have sufficed for this full-group conclusion.

For finite connected unit-conductance networks, resistance is the minimum
of
\[
 \energy(\theta)=\sum_{\{x,y\}}\theta(x,y)^2
\]
over unit flows between the specified terminals; sum once per
\emph{unoriented} edge. We use Thomson's principle, Rayleigh monotonicity
and Kirchhoff's edge-marginal identity \cite[\S4, Theorem~4.1]{BLPS}.
Define free resistance by a finite free induced exhaustion:
\[
 \RF_X(x,y)=\lim_n R_{X_n}(x,y).
\]
The limit decreases and is exhaustion-independent. The finite identity
and the free weak limit give, for every edge \(e=\{x,y\}\),
\begin{equation}\label{marginal}
 \mathbb P(e\in\FUSF(X))=\RF_X(x,y).
\end{equation}
See also \cite[Proposition~7.1 and Theorem~7.8]{BLPS}.
Connected edge-exhaustions give the same limit: they are eventually
sandwiched between any fixed finite induced ball and some larger ball,
so Rayleigh monotonicity squeezes the resistance limits.
These facts concern free boundaries, not wired resistance or the
unconstrained minimum among infinite square-summable flows.

\section{First calculation: branches and a diagonal free exhaustion}
For a vertex \(u\), write \(a=p(u)\) and \(b=p^2(u)\).
Let \(B_m(u;a,b)\) contain \(u\), its descendants through \(m\) generations,
and terminals \(a,b\). Include every grandparent-graph edge incident to
an interior vertex with both endpoints in this set, and omit the edge
\(ab\). An untruncated descendant branch meets the rest of \(G\) only at
\(a,b\). The truncation deletes external edges; it makes no identifications.
Different child branches have disjoint interiors.

Let \(\kappa_m\) be the terminal conductance of \(B_m\).
The case \(m=0\) is the two-unit-edge path \(a-u-b\).
Replacing each of \(u\)'s two child branches by its terminal conductance
therefore gives the exact recurrence
\begin{equation}\label{krec}
 \kappa_0=\tfrac12,\qquad
 \kappa_{m+1}=\frac{1+2\kappa_m}{2+2\kappa_m}.
\end{equation}
Indeed the \(u a\) conductance becomes \(1+2\kappa_m\), in series
with the unit edge \(u b\). The map in \eqref{krec} is increasing,
starts the sequence increasing from \(1/2\), and is bounded above by \(1\).
Its unique positive fixed point is
\[
 \kappa=1/\sqrt2,\qquad A=1+2\kappa=1+\sqrt2,\qquad
 \kappa=\frac{A}{1+A}.
\]

Put \(x_i=p^i(o)\).
For \(m\ge0,n\ge2\), form the following finite connected subgraph \(L_{m,n}\):
the spine \(x_0,\ldots,x_n\) with all adjacent and skip-one edges;
both depth-\(m\) child branches at \(x_0\), with terminals \(x_0,x_1\);
and the one depth-\(m\) off-spine child branch at each \(x_i\),
\(1\le i<n\), with terminals \(x_i,x_{i+1}\).
All original distinct edges retain conductance one.
These graphs increase with each parameter.
Every end-ray from a fixed vertex merges with the spine at a finite place;
its off-spine depth is finite. Thus \(L_{n,n}\) eventually contains every
vertex and every edge of \(G\). The temporarily omitted branches at the
top \(x_n\) appear at later stages. This is a genuine free edge-exhaustion.

Eliminate the attached branches. The conductance on \(x_0x_1\) is
\(A_m=1+2\kappa_m\); on each higher adjacent pair it is
\(C_m=1+\kappa_m\); each skip-one edge still has conductance one.
For the upper suffix \(x_1,\ldots,x_n\), its terminal conductance is
\begin{equation}\label{upper}
 B_{m,n}=T_{C_m}^{\,n-2}(C_m),\qquad
 T_C(t)=C+\frac{t}{1+t}.
\end{equation}
This follows by eliminating the top triangular cell successively:
a direct \(C\) is parallel to a unit edge in series with the remaining
upper conductance.

For \(t\ge C_m\ge1\), \(T'_{C_m}(t)\le1/4\).
Here \(C_m\) stays in a compact interval and the positive fixed point
of \(T_C\), the positive root of \(B^2-CB-C=0\), depends continuously on
\(C\). Hence the convergence of the iterates in \eqref{upper} is
geometrically uniform in \(m\).
As \(m=n\to\infty\), \(C_m\to C=1+\kappa\) and \(B_{n,n}\to A\),
since \(A=C+A/(1+A)\). This justifies the actual diagonal limit;
no unsupported interchange of two infinite limits is made.

The reduced three-terminal network has conductances \(A_m,B_{m,n},1\)
on \(x_0x_1,x_1x_2,x_0x_2\). Thus
\begin{align}
 R_{L_{m,n}}(x_0,x_1)&=
 \left(A_m+\frac{B_{m,n}}{1+B_{m,n}}\right)^{-1},\label{tri1}\\
 R_{L_{m,n}}(x_0,x_2)&=
 \left(1+\frac{A_mB_{m,n}}{A_m+B_{m,n}}\right)^{-1}.\label{tri2}
\end{align}
Along the free exhaustion, these converge respectively to
\begin{equation}\label{res}
 R_{\rm tree}=\frac1{1+3\kappa}=\frac{3\sqrt2-2}{7},
 \qquad
 R_{\rm gp}=\frac2{3+2\kappa}=\frac{6-2\sqrt2}{7}.
\end{equation}
The end stabilizer is transitive on unoriented edges of each type.
Applying \eqref{marginal} separately to the three tree edges and five
grandparent edges at \(o\) gives
\[
 \E\deg_{\FUSF(G)}(o)=3R_{\rm tree}+5R_{\rm gp}
                  =\frac{24-\sqrt2}{7}.
\]
Since \((73/50)^2>2\), the last expression exceeds
\((24-73/50)/7=161/50\). No mass-transport argument is used.

\section{Second calculation: an induced exhaustion and contraction}
We independently recover \eqref{res} using only induced finite graphs.
Put \(v_j=p^j(o)\), and let \(S_n\) be all descendants of \(v_n\)
through depth \(2n\). Let \(G_n=G[S_n]\), \(n\ge2\).
When \(n\) increases by one, old depths increase by one but the allowed
depth increases by two. These finite connected sets are nested and
exhaust \(T\), by end-ray coalescence.

Consider descendants of \(x\) through depth \(m\), together with its
parent \(y\), retaining the direct edge \(xy\) but excluding the edge
from \(x\) to its grandparent. Eliminating the strict descendants gives
conductance \(c_m\) between \(x,y\), where
\begin{equation}\label{crec}
 c_0=1,\qquad c_m=1+\frac{2c_{m-1}}{1+c_{m-1}},
 \qquad a_m=1+\frac{c_{m-1}}{1+c_{m-1}}\quad(m\ge1).
\end{equation}
Each of the two children supplies its reduced \(c_{m-1}\) in series
with a unit edge to \(y\), proving the recurrence directly.

Eliminate off-spine vertices in \(G_n\).
The resulting conductances on the spine \(v_0,\ldots,v_n\) are
\[
 v_0v_1:c_n,\qquad v_jv_{j+1}:a_{n+j}\ (1\le j<n),
 \qquad v_jv_{j+2}:1\ (0\le j\le n-2).
\]
Both children at \(v_0\) are off-spine; above it there is one off-spine
child, whose descendant depth is \(n+j-1\).
The off-spine part at the top \(v_n\) has only one attachment and has no
effect; no terminal outside \(G_n\) is introduced.
Reducing the upper suffix gives
\begin{equation}\label{urec}
 u_{n-1,n}=a_{2n-1},\qquad
 u_{j,n}=a_{n+j}+\frac{u_{j+1,n}}{1+u_{j+1,n}}
       \quad(1\le j\le n-2).
\end{equation}
The final triangle has conductances \(c_n,u_{1,n},1\).

The map \(t\mapsto1+2t/(1+t)\) contracts on \([1,\infty)\) by at most
\(1/2\), so \(c_m\to c=1+\sqrt2=A\) and
\(a_m\to a=1+c/(1+c)\).
For the suffix maps in \eqref{urec} the derivative in their argument
is at most \(1/4\), all arguments and coefficients are uniformly bounded,
and every coefficient index is at least \(n+1\).
With \(\epsilon_n=\sup_{m\ge n+1}|a_m-a|\), iteration gives
\[
 |u_{1,n}-c|
 \le \frac43\epsilon_n+C_0\,4^{-(n-2)}\longrightarrow0,
\]
because \(c=a+c/(1+c)\). Apply \eqref{tri1}--\eqref{tri2} with
\(c_n,u_{1,n}\) and use this induced free exhaustion.
The limits are again exactly \eqref{res}.

For reproducible finite checks, global rational elimination on the actual
induced networks gives, at \(n=2\), the triangle
\((7/3,17/10,1)\) and resistances \(27/80,121/240\).
At \(n=3\) it gives \((12/5,6236/2669,1)\).
These checks support the finite indexing; the written contraction proofs
establish the all-depth and infinite-volume assertions.

\section{A complete finite certificate for the square}
Let \(W\) be the radius-\(10\) ball about \(o\) in \(T\).
Write \(K=G[W]\), \(J=K^2\). Then \(J\subseteq H\).
Every \(H\)-neighbor of \(o\) has a witness using at most two \(G\)-edges,
hence a tree walk of at most four edges, entirely within the radius-\(4\)
tree ball. Therefore \(J\) contains \emph{every} \(H\)-edge at \(o\).

For clarity, the complete neighbor count is as follows:
four strict ancestors; \(2+4+8+16=30\) strict descendants in depths
one through four; the sibling; the uncle (the parent's sibling);
the sibling's two children; and the uncle's two children.
These disjoint possibilities exhaust two ancestor-or-descendant \(G\)
steps and total \(40\). The end stabilizer preserves \(H\), so its degree
is \(40\) everywhere.

Here is an exact specification of the finite data.
Index the tree ball by breadth-first lexicographic word order:
the root has index \(0\) and children \(1,2,3\); vertex \(i\ge1\)
has outward children \(2i+2,2i+3\).
Thus
\[
 |W|=1+3(2^{10}-1)=3070,\qquad |E(J)|=19931.
\]
Choose the distinguished end ray \(0,1,4,10,22,\ldots\).
On this ray \(p(0)=1\), \(p(i)=2i+2\) for \(i\ne0\);
off this ray \(p(i)\) is the outward-tree predecessor
(\(0\) for \(i=1,2,3\), otherwise \(\lfloor(i-2)/2\rfloor\)).
Construct \(K\) by all pairs \(i,p(i)\) and \(i,p^2(i)\) with endpoints
in \(W\), then square \(K\).
Order the unordered \(J\)-edges \((i,j)\), \(i<j\), lexicographically.
The sorted root-neighbor indices are
\[
 1,\ldots,22,\quad24,25,\quad30,\ldots,45.
\]

The source supplement contains the exact compressed JSON certificate
named square\_\allowbreak{}flow\_\allowbreak{}certificate.json.gz, and a short independent
standard-library verifier.
For each of those \(40\) targets \(v\), its row supplies one integer
\(n_{v,e}\) for every ordered-list edge \(e\).
Give that edge its listed orientation \(i\to j\), and put
\(\theta_v(e)=n_{v,e}/D\), \(D=10^9\).
The verifier regenerates the \emph{whole} finite graph from the indexed
tree, matches every edge and terminal, and checks the following finite
integer identities:
\begin{align}
 \sum_{e\text{ out of }x}n_{v,e}
 -\sum_{e\text{ into }x}n_{v,e}
 &=D\bigl(\mathbf1_{\{x=v\}}-\mathbf1_{\{x=o\}}\bigr),
       &&x\in W,\label{div}\\
 \sum_v\sum_{e\in E(J)}n_{v,e}^2
 &=3208526010913748780.\label{cost}
\end{align}
All \(797240\) current entries are actual integers;
all \(122800\) vertex-divergence equations and all \(40\) individual
energy numerators are verified. Each unoriented edge is counted once.
This is feasible-current verification, not trust in a numerical residual
or a claim of current optimality. Floating-point solves were used to
find the original witnesses; their generator is unnecessary for checking
\eqref{div}--\eqref{cost}. The immutable data and exact checking suffice.

\begin{lemma}\label{transfer}
For the currents above,
\[
 \E\deg_{\FUSF(H)}(o)
 \le\sum_v\energy(\theta_v)
 =\frac{3208526010913748780}{10^{18}}<\frac{321}{100}.
\]
\end{lemma}
\begin{proof}
Equation~\eqref{div} makes \(\theta_v\) a unit flow from \(v\) to \(o\)
in the finite graph \(J\). Thomson's principle gives
\(R_J(o,v)\le\energy(\theta_v)\).
Take any connected finite induced free exhaustion \(H_n\) of \(H\).
For large \(n\), \(J\subseteq H_n\), so extend \(\theta_v\) by zero.
It remains a unit flow of unchanged energy, giving
\[
 R_{H_n}(o,v)\le R_J(o,v)\le\energy(\theta_v).
\]
Pass to the free limit and apply \eqref{marginal}.
Summing over the \emph{complete} root neighborhood gives the claim
from \eqref{cost}. The final strict inequality is the integer comparison
\(100\cdot3208526010913748780<321\cdot10^{18}\).
\end{proof}

Only finitely supported flows were transferred to finite free networks.
Minimizing over arbitrary square-summable infinite flows instead would
give wired resistance and would not justify the required free upper
bound. No boundary is shorted and no wired bound is substituted.
Different root edges may use different flows, since only marginals are
summed.

\begin{proof}[Completion of Theorem~\ref{main}]
The graph assertions and \eqref{gvalue} were proved above.
Lemma~\ref{transfer} proves \eqref{hbound}. With \eqref{bf},
\[
 b_F(G)-b_F(H)
 >\frac12\left(\frac{161}{50}-\frac{321}{100}\right)=\frac1{200}.
\]
\end{proof}

\section{Predecessors, scope, and reproducibility}
The grandparent graph, finite electrical principles and free-exhaustion
framework are inherited, with explicit references above.
Tang's inspected arXiv version proves connectivity and branching
properties for this family, not the comparison proved here.
Pete--Tim\'ar \cite[Theorems~1.1--1.2]{PT} already show that FUSF
\emph{component counts} can depend on the generating set even for a
fixed virtually free group; they explicitly distinguish this from
expected degree, which is group-invariant for Cayley graphs.
Those results are not claimed as new here.
Gaboriau's harmonic-projection invariant is formulated with a
transitive unimodular action \cite[\S2.4]{Gab}.
No such tracial hypothesis is used in our negative calculation.
No claim concerning failure in the unimodular subclass is made.

The original question was checked in the actual seven-page 2010 source,
including its unrestricted scope and Cayley comment.
Our bounded primary-source investigation through 22 October 2026
authenticated no earlier publication of this complete expected-degree
square comparison. That does not certify historical firstness or
continued worldwide openness.
Tang's final journal PDF and Pete--Tim\'ar's final journal PDF returned
access-challenge pages; their actual arXiv versions identified below
were read instead. The full Lyons--Peres book and all unindexed sources
were not inspected. Source-edition and exact read-scope limits are
recorded in the supplement. None of these dependency DOIs is presented
as a prior publication of this full counterexample.

The supplement contains the entire rational-current data, the independent
indexed-tree verifier, exact per-target outcomes, an independent
binary-word/global-rational-elimination check of the finite reductions,
editable manuscript, dependencies, priority notes, licenses and
offline packaging instructions.
Every current certificate entry is checked exactly; no external solver,
numerical tolerance, original claimant program or installed third-party
package is needed for replay.
The certificate proves a finite statement; Lemma~\ref{transfer} is the
separate universal argument making it an infinite free-forest bound.
The finite circuit diagnostics do not prove the infinite limits;
Sections~3--4 do. No formal proof-assistant certification is claimed.

\paragraph{Disclosure and license.}
AI tools were extensively involved in mathematical research and
exploration, proof development and reconstruction, computational
verification, literature and source investigation, manuscript preparation,
editing, and adversarial auditing.
Alec Kriebel is the named human author.
This is an unrefereed preprint and has not undergone formal independent
human peer review. AI reviews do not supply human refereeing or journal
acceptance. Original prose and certificate data are CC BY 4.0;
original verification and packaging code is MIT licensed.

\begin{thebibliography}{9}
\raggedright
\bibitem{Abert}
M.~Ab\'ert, \emph{Some questions}, 2 November 2010,
printed p.~4, Question~24.
\href{https://www.renyi.hu/~abert/questions.pdf}{Author-hosted problem list}.
\bibitem{BLPS}
I.~Benjamini, R.~Lyons, Y.~Peres, and O.~Schramm,
\emph{Uniform spanning forests},
Ann.~Probab. \textbf{29} (2001), 1--65.
\href{https://doi.org/10.1214/aop/1008956321}{doi:10.1214/aop/1008956321}.
Inspected corrected author version dated 27 January 2009:
\href{https://rdlyons.pages.iu.edu/pdf/usf.pdf}{usf.pdf}.
\bibitem{Gab}
D.~Gaboriau, \emph{Invariant percolation and harmonic Dirichlet functions},
Geom.~Funct.~Anal. \textbf{15} (2005), 1004--1051.
\href{https://doi.org/10.1007/s00039-005-0539-2}{doi:10.1007/s00039-005-0539-2}.
Inspected arXiv:math/0405458v3.
\bibitem{PT}
G.~Pete and \'A.~Tim\'ar,
\emph{The free uniform spanning forest is disconnected in some virtually
free groups, depending on the generator set},
Ann.~Probab. \textbf{50} (2022), 2218--2243.
\href{https://doi.org/10.1214/22-AOP1581}{doi:10.1214/22-AOP1581}.
Actual inspected body: arXiv:2006.06387v2, 24 January 2021.
\bibitem{Tang}
P.~Tang,
\emph{Weights of uniform spanning forests on nonunimodular transitive graphs},
Electron.~J.~Probab. \textbf{26} (2021), 1--62.
\href{https://doi.org/10.1214/21-EJP709}{doi:10.1214/21-EJP709}.
Actual inspected body: arXiv:1908.09889v2, 3 December 2020.
\end{thebibliography}
\end{document}
